\documentclass[11pt,a4paper]{article}

\usepackage{amsmath,amssymb,amsthm}
\usepackage{mathtools}
\usepackage{bm}
\usepackage{graphicx}
\usepackage{hyperref}
\usepackage{booktabs}
\usepackage{xcolor}
\usepackage[margin=1.1in]{geometry}
\usepackage{algorithm}
\usepackage{algpseudocode}
\usepackage{natbib}

% Theorem environments
\newtheorem{theorem}{Theorem}[section]
\newtheorem{proposition}[theorem]{Proposition}
\newtheorem{lemma}[theorem]{Lemma}
\newtheorem{corollary}[theorem]{Corollary}
\newtheorem{definition}[theorem]{Definition}
\newtheorem{remark}[theorem]{Remark}

% Notation shortcuts
\newcommand{\R}{\mathbb{R}}
\newcommand{\E}{\mathbb{E}}
\newcommand{\calL}{\mathcal{L}}
\newcommand{\calZ}{\mathcal{Z}}
\newcommand{\calX}{\mathcal{X}}
\newcommand{\calT}{\mathcal{T}}
\newcommand{\calB}{\mathcal{B}}
\newcommand{\norm}[1]{\left\|#1\right\|}
\newcommand{\Cov}{\operatorname{Cov}}
\newcommand{\tr}{\operatorname{tr}}
\newcommand{\diag}{\operatorname{diag}}
\newcommand{\rank}{\operatorname{rank}}
\newcommand{\Var}{\operatorname{Var}}

\title{\textbf{Linear Disentanglement in Representation Learning:\\
Definition, Projection, and Knowledge Distillation}}

\author{Technical Report}
\date{}

\begin{document}
\maketitle

\begin{abstract}
Recent self-supervised vision models, notably DINO~\citep{caron2021emerging} and
DINOv2~\citep{oquab2023dinov2}, produce features with a remarkable property:
they are \emph{linearly separable} — a simple linear probe achieves near-supervised
accuracy on ImageNet without any fine-tuning.  Despite widespread empirical
observation of this property, a precise mathematical characterisation and a
principled method to \emph{transfer} it to a student model via knowledge
distillation remain underexplored.
This report provides three contributions.
First, we formalise \emph{linear disentanglement} as a structural property of a
feature distribution, decomposing it into two independently measurable axes:
\emph{isotropy} (eigenvalue uniformity of the covariance) and
\emph{decorrelation} (vanishing off-diagonal covariance).
Second, we show that ZCA whitening is the canonical projection into the linearly
disentangled space and characterise its optimality among orthogonal
transformations.
Third, we derive a three-term distillation objective — variance regularisation,
covariance regularisation, and whitened-space alignment — that simultaneously
\emph{enforces} the structural property on the student and \emph{transfers}
semantic content from the teacher.
We discuss connections to VICReg~\citep{bardes2022vicreg}, Barlow
Twins~\citep{zbontar2021barlow}, and Fisher Linear Discriminant Analysis.
\end{abstract}

\tableofcontents
\newpage

%%% ─────────────────────────────────────────────────────────────────────────
\section{Introduction}
\label{sec:intro}
%%% ─────────────────────────────────────────────────────────────────────────

Knowledge distillation typically trains a student network to mimic the
\emph{output} of a teacher network — matching class probabilities, intermediate
activations, or pairwise relational structure.
Such objectives transfer \emph{what} the teacher knows but not necessarily
\emph{how} the teacher organises its internal representations.

DINO-style teachers are special: their patch tokens form rich spatial maps that
can be directly clustered into semantically coherent segments without any
supervision~\citep{caron2021emerging}.
This is not merely a consequence of high accuracy; it follows from a structural
property of the feature distribution.
Concretely, DINO features are approximately \emph{isotropic} and
\emph{decorrelated}, meaning the representation space is efficiently tiled by
semantic directions with roughly equal energy.
This property — which we formalise as \emph{linear disentanglement} — is what
makes linear probes succeed.

A student trained only to minimise cosine distance to teacher features will
reproduce the directions of teacher features but may not reproduce their
structural organisation.
In particular, standard cosine or MSE losses are invariant to isotropic scaling
and rotation: they do not penalise a student for concentrating all variance into
a few directions (anisotropy) or for allowing correlated dimensions
(entanglement).

This report addresses the question: \textit{how can we distill not only the
semantic content of a teacher but also its linear disentanglement property?}

%%% ─────────────────────────────────────────────────────────────────────────
\section{Notation and Setup}
\label{sec:notation}
%%% ─────────────────────────────────────────────────────────────────────────

Let $f_\theta : \calX \to \R^D$ be a student vision encoder parameterised by
$\theta$.  For a minibatch of $N$ images $\{x_i\}_{i=1}^N$, define the feature
matrix
\[
    Z \;=\; [f_\theta(x_1),\, \ldots,\, f_\theta(x_N)]^\top \;\in\; \R^{N \times D}.
\]
For a ViT backbone, each image yields $L$ patch tokens; we write
$Z \in \R^{B L \times D}$ when the batch and spatial dimensions are flattened.

The centred feature matrix is $\tilde{Z} = Z - \mathbf{1}\bar{z}^\top$ where
$\bar{z} = \frac{1}{N}\sum_i z_i$.  The (unbiased) sample covariance is
\[
    C \;=\; \frac{1}{N-1}\,\tilde{Z}^\top \tilde{Z} \;\in\; \R^{D \times D}.
\]
Let $C = U \Sigma U^\top$ be the eigendecomposition with
$\Sigma = \diag(\sigma_1 \ge \sigma_2 \ge \cdots \ge \sigma_D \ge 0)$.

A teacher encoder $g_\phi : \calX \to \R^{D_t}$ is fixed and produces teacher
features $Z^{(t)} \in \R^{N \times D_t}$.  A projection head
$h_\psi : \R^D \to \R^{D_t}$ maps student features into the teacher's dimension.


%%% ─────────────────────────────────────────────────────────────────────────
\section{Defining Linear Disentanglement}
\label{sec:definition}
%%% ─────────────────────────────────────────────────────────────────────────

\subsection{Classical Linear Separability vs.\ Representational Disentanglement}
\label{sec:classical}

The classical notion of linear separability states that a dataset
$\{(z_i, y_i)\}$ is \emph{linearly separable} if there exists a weight vector
$w \in \R^D$ and bias $b$ such that $w^\top z_i + b > 0$ whenever $y_i = +1$ and
$< 0$ otherwise.
This is a binary, label-dependent property.

For representation learning, we require a richer, label-free notion that
characterises the \emph{geometry} of the feature distribution.
Intuitively, a representation is linearly disentangled if

\begin{enumerate}
    \item it uses each dimension independently to encode distinct factors of
          variation (decorrelation), and
    \item it allocates roughly equal energy to each dimension rather than
          collapsing into a low-dimensional subspace (isotropy).
\end{enumerate}

\noindent These conditions do not by themselves guarantee that classes are
separable — they constrain the unlabelled marginal distribution, not the class
geometry (Section~\ref{sec:joint} makes this precise).  Rather, they make a
linear probe \emph{efficient to fit}: (i) no dimension is redundant and (ii) no
single direction dominates the gradient signal during probing, so the probe is
well-conditioned and label-efficient whenever the classes are in fact separable.

\subsection{Isotropy: Eigenvalue Uniformity}
\label{sec:isotropy}

\begin{definition}[Participation Ratio]
\label{def:pr}
The \emph{participation ratio} (PR) of a feature matrix $Z \in \R^{N \times D}$
is
\[
    \rho(Z) \;=\; \frac{\bigl(\sum_{i=1}^r \sigma_i\bigr)^2}{\,r \sum_{i=1}^r \sigma_i^2\,}
    \;\in\; \bigl(0,\, 1\bigr],
\]
where the $\sigma_i$ are the eigenvalues of the sample covariance $C$
(equivalently $\sigma_i \propto s_i^2$ for the singular values $s_i$ of
$\tilde{Z}$) and $r = \rank(C)$.  Throughout, all spectral diagnostics
($\rho$, spectral entropy, and the effective rank of
Section~\ref{sec:topk}) are defined on this single eigenvalue (variance)
spectrum, so they are mutually consistent.
\end{definition}

\begin{proposition}
$\rho(Z) = 1$ if and only if $\sigma_1 = \sigma_2 = \cdots = \sigma_r$
(perfect isotropy). $\rho(Z) = 1/r$ if all variance is concentrated in a single
direction.
\end{proposition}
\begin{proof}
By the Cauchy-Schwarz inequality,
$\bigl(\sum_i \sigma_i\bigr)^2 \le r \sum_i \sigma_i^2$
with equality iff all $\sigma_i$ are equal.
Dividing by $r \sum_i \sigma_i^2$ gives $\rho \le 1$ with $\rho = 1$ at equality.
For the lower bound, when $\sigma_1 > 0$ and $\sigma_i = 0$ for $i > 1$,
$\rho = \sigma_1^2 / (r \sigma_1^2) = 1/r$.
\end{proof}

A related measure is the \emph{spectral entropy}:
\[
    H(Z) \;=\; -\sum_{i=1}^r p_i \log p_i, \quad
    p_i = \frac{\sigma_i}{\sum_j \sigma_j}.
\]
This is maximised at $H = \log r$ when all singular values are equal, recovering
the uniform (isotropic) case.  Both $\rho$ and $H$ are label-free diagnostics
for collapse detection during training.

\subsection{Decorrelation: Off-Diagonal Covariance}
\label{sec:decorrelation}

\begin{definition}[Decorrelation Penalty]
\label{def:decor}
The \emph{decorrelation penalty} of feature matrix $Z$ is
\[
    \mathcal{D}(Z)
    \;=\; \norm{C \odot (\bm{J} - \bm{I})}_F^2
    \;=\; \sum_{i \ne j} C_{ij}^2,
\]
where $\bm{J}$ is the all-ones matrix, $\odot$ is element-wise product, and
$C$ is the sample covariance.
\end{definition}

$\mathcal{D}(Z) = 0$ if and only if the feature dimensions are pairwise
uncorrelated, i.e.\ $C$ is diagonal.  Combined with Definition~\ref{def:pr},
full linear disentanglement corresponds to $C = \sigma^2 I$ for some
$\sigma > 0$.

\begin{remark}[Information-Theoretic Interpretation]
Minimising $\mathcal{D}(Z)$ is a second-order approximation to minimising the
pairwise mutual information $\sum_{i \ne j} I(Z_i; Z_j)$ between feature
dimensions.  Under Gaussianity, the approximation is exact: $I(Z_i; Z_j) = 0
\Leftrightarrow C_{ij} = 0$.
\end{remark}

\subsection{Joint Characterisation}
\label{sec:joint}

\begin{definition}[Linear Disentanglement]
\label{def:ld}
A feature distribution $p(z)$ with mean $\mu$ and covariance $C$ is
\emph{$(\rho_0, \varepsilon)$-linearly disentangled} if
\begin{equation}
    \rho(Z) \;\ge\; \rho_0, \qquad
    \mathcal{D}(Z) \;\le\; \varepsilon.
\label{eq:ld_def}
\end{equation}
In the ideal limit $(\rho_0, \varepsilon) = (1, 0)$, the covariance reduces to a
scaled identity: $C = \sigma^2 I$.
\end{definition}

\noindent\textbf{Connection to Linear Probe Performance.}\;
We stress that linear disentanglement is a property of the \emph{marginal}
(total) feature covariance $C$, which decomposes as $C = S_W + S_B$ into
within- and between-class scatter $S_W, S_B$.  Consequently $C = \sigma^2 I$
does \emph{not} by itself imply class separability: it constrains only the sum
$S_W + S_B$, not the class geometry, and one can have $C = \sigma^2 I$ with
coincident class means (zero separability).  Disentanglement should therefore be
read as a \emph{conditioning} property rather than a sufficient condition for a
successful probe.  Two precise statements make the connection rigorous.

\emph{(i) Whitening turns the optimal probe into nearest-class-mean.}
Under the homoscedastic-Gaussian model $z \mid y{=}c \sim
\mathcal{N}(\mu_c, S_W)$ with shared within-class covariance $S_W$, the
Bayes-optimal classifier is linear and ranks classes by the Mahalanobis distance
$\norm{z-\mu_c}_{S_W^{-1}}$.  Applying the $S_W$-whitening $z \mapsto
S_W^{-1/2}z$ makes this distance Euclidean, so the optimal probe reduces to
nearest-class-mean and separability is governed by the \emph{whitened}
between-class separation $\norm{S_W^{-1/2}(\mu_a-\mu_b)}_2$.  This is the precise
sense in which whitening is the best case for a linear probe — but it is
whitening by $S_W$, not by the total covariance $C$, and it recovers exactly
Fisher's criterion $J = \tr(S_W^{-1}S_B)$.

\emph{(ii) Isotropy improves probe conditioning.}
When the probe is fit by (regularised) least squares or gradient descent, its
optimisation is governed by the condition number $\kappa(C) =
\sigma_1/\sigma_D$.  Isotropy ($\rho \to 1$, i.e. $\kappa \to 1$) minimises this
condition number, giving faster, more label-efficient probe convergence and
tighter margin-based generalisation bounds, \emph{independently} of the class
geometry.

The objective of Section~\ref{sec:distillation} whitens by the total batch
covariance — an unsupervised proxy that coincides with $S_W$ only when
$S_B \ll S_W$.  It therefore delivers (ii) directly and approximates (i); a
fully supervised $S_W$-whitening would be required to realise Fisher's
discriminant exactly.

%%% ─────────────────────────────────────────────────────────────────────────
\section{Transforming Features into the Linearly Disentangled Space}
\label{sec:whitening}
%%% ─────────────────────────────────────────────────────────────────────────

\subsection{ZCA Whitening}
\label{sec:zca}

Given features $Z \in \R^{N \times D}$, the goal is to find a linear
transformation $W \in \R^{D \times D}$ such that the transformed features
$Z_w = \tilde{Z} W$ satisfy $\Cov(Z_w) = I$.

\begin{definition}[ZCA Whitening Transform]
\label{def:zca}
Let $C = U \Sigma U^\top$ be the eigendecomposition of the sample covariance.
The \emph{ZCA whitening matrix} is
\begin{equation}
    W_{\mathrm{ZCA}} \;=\; U \,\Sigma^{-1/2} U^\top,
\label{eq:zca}
\end{equation}
where $\Sigma^{-1/2} = \diag(\sigma_1^{-1/2}, \ldots, \sigma_D^{-1/2})$
(with regularisation $\sigma_i \leftarrow \sigma_i + \varepsilon$ in practice).
The whitened features are $Z_w = \tilde{Z}\, W_{\mathrm{ZCA}}$.
\end{definition}

\begin{proposition}[Covariance of Whitened Features]
$\Cov(Z_w) = I_D$.
\end{proposition}
\begin{proof}
\[
    \Cov(Z_w) = W_{\mathrm{ZCA}}^\top C\, W_{\mathrm{ZCA}}
    = (U \Sigma^{-1/2} U^\top)^\top (U \Sigma U^\top) (U \Sigma^{-1/2} U^\top)
    = U \Sigma^{-1/2} \Sigma \Sigma^{-1/2} U^\top = I.
\]
\end{proof}

\noindent After ZCA whitening, the feature distribution achieves perfect linear
disentanglement in the sense of Definition~\ref{def:ld} with
$(\rho_0, \varepsilon) = (1, 0)$.  We caution that this is a statement about the
marginal covariance only: when the spectrum is highly anisotropic, exact
full-rank whitening amplifies near-zero directions and can \emph{degrade} rather
than improve probe accuracy.  The regularised, rank-truncated variant of
Section~\ref{sec:topk} is therefore what we use in practice; full-rank ZCA is the
idealised limit, not the deployed transform.

\subsection{Optimality Among Whitening Transforms}
\label{sec:zca_optimality}

Many whitening transforms satisfy $\Cov(Z_w) = I$.  The PCA whitening matrix
$W_{\mathrm{PCA}} = U \Sigma^{-1/2}$ also achieves this but additionally rotates
features into the eigenbasis of $C$.  ZCA whitening has a unique optimality
property:

\begin{theorem}[ZCA Minimises Rotation]
\label{thm:zca_optimal}
Among all whitening matrices $W$ satisfying $W^\top C W = I$, the ZCA
transform $W_{\mathrm{ZCA}}$ minimises the expected squared distance between
original and whitened features:
\[
    W_{\mathrm{ZCA}} = \arg\min_{W:\,W^\top C W = I}\; \E_z \norm{z - Wz}_2^2.
\]
\end{theorem}
\begin{proof}
We write the per-sample whitening as $z \mapsto W^\top z$ (consistent with
$Z_w = \tilde Z\,W$); since $W_{\mathrm{ZCA}}$ is symmetric this coincides with
$z \mapsto Wz$ at the optimum.  With $z$ centred and $\E[zz^\top] = C$,
\[
    \E\norm{z - W^\top z}^2
    = \tr\!\bigl((I - W)^\top C\,(I - W)\bigr)
    = \tr(C) - 2\tr(CW) + \tr(W^\top C W).
\]
Every whitening matrix satisfies $W^\top C W = I$, so $\tr(W^\top C W) = D$ is
constant and minimising the distance is equivalent to \emph{maximising}
$\tr(CW)$.  The constraint $W^\top C W = I$ is equivalent to $W = C^{-1/2}Q$ for
some orthogonal $Q$, where $C^{-1/2} = U\Sigma^{-1/2}U^\top$ is the symmetric
inverse square root (indeed
$W^\top C W = Q^\top C^{-1/2} C\, C^{-1/2} Q = Q^\top Q = I$).  Then
\[
    \tr(CW) = \tr\!\bigl(C\,C^{-1/2}Q\bigr) = \tr\!\bigl(C^{1/2}Q\bigr).
\]
For the symmetric positive-definite matrix $C^{1/2}$, the orthogonal Procrustes
inequality gives $\tr(C^{1/2}Q) \le \tr(C^{1/2})$ over all orthogonal $Q$, with
equality iff $Q = I$.  Hence the maximiser is
$W = C^{-1/2} = U\Sigma^{-1/2}U^\top = W_{\mathrm{ZCA}}$.
\end{proof}

This optimality has a direct implication for distillation: applying ZCA
whitening to teacher features produces a target that is as close as possible
(in $\ell_2$ sense) to the original teacher features while being perfectly
linearly disentangled.  The whitening does not arbitrarily rotate the semantic
directions; it merely rescales them.

\subsection{Batch Approximation and Regularisation}
\label{sec:batch_zca}

During training, the true covariance $C$ is unknown and must be estimated from
minibatches of size $N$.  The estimator
\[
    \hat{C} = \frac{1}{N}\, \tilde{Z}^\top \tilde{Z}
\]
is consistent but biased for small $N$.  In practice two adjustments are needed.

\textbf{Regularisation.}\; Near-zero eigenvalues (dormant dimensions) cause
numerical instability in $\Sigma^{-1/2}$.  We replace $\sigma_i$ by
$\sigma_i + \varepsilon$ with $\varepsilon \in [10^{-4}, 10^{-2}]$.
This also controls the \emph{whitening radius}: strongly regularised whitening
$({\varepsilon} \gg \sigma_i)$ approaches an isotropic identity transform, while
$\varepsilon \to 0$ recovers exact ZCA.

\textbf{Feature dimension.}\; For high-dimensional features ($D \gg N$), the
empirical covariance has rank at most $N$, and the pseudo-inverse should be used.
In practice, for ViT patch features with $D = 768$ and $N = B \cdot L \ge 4096$,
the batch is typically large enough for a well-conditioned estimate.

\textbf{Computational cost.}\; The dominant cost is the EVD (or SVD) of
$\hat{C} \in \R^{D \times D}$, which is $\mathcal{O}(D^3)$ per step.  For
$D = 768$, this is approximately $4.5 \times 10^8$ FLOPs — comparable to one
transformer attention layer.  The whitening matrix $W$ is applied as a
single matmul $\tilde{Z} W$ at cost $\mathcal{O}(ND^2)$.

%%% ─────────────────────────────────────────────────────────────────────────
\section{Knowledge Distillation of Linear Disentanglement}
\label{sec:distillation}
%%% ─────────────────────────────────────────────────────────────────────────

\subsection{Problem Formulation}
\label{sec:formulation}

Standard feature-level distillation minimises
\[
    \calL_{\mathrm{distill}}
    = \E_{x}\Bigl[\ell\bigl(h_\psi(f_\theta(x)),\; g_\phi(x)\bigr)\Bigr],
\]
where $\ell$ is typically cosine distance or MSE.
This objective is invariant to arbitrary invertible linear transformations of
the teacher's feature space, meaning the student may converge to a
representation that has the same pairwise angles as the teacher but completely
different covariance structure.

Our goal is to design $\calL$ such that the student satisfies
Definition~\ref{def:ld} at convergence \emph{without requiring class labels}.

\subsection{The Three-Term Objective}
\label{sec:loss}

Let $Z_s = h_\psi(f_\theta(\mathcal{B})) \in \R^{N \times D_t}$ be student
features after the projection head on a minibatch $\mathcal{B}$, and
$Z_t = g_\phi(\mathcal{B}) \in \R^{N \times D_t}$ be teacher features
(stop-gradient).  We propose:

\begin{equation}
\boxed{
    \calL_{\mathrm{LD}}(Z_s, Z_t)
    \;=\; \lambda_v \calL_v(Z_s)
    \;+\; \lambda_c \calL_c(Z_s)
    \;+\; \lambda_a \calL_a(Z_s, Z_t)
}
\label{eq:total_loss}
\end{equation}

with the three terms defined as follows.

\paragraph{Variance term $\calL_v$ — prevent dimensional collapse.}
\begin{equation}
    \calL_v(Z_s) \;=\; \frac{1}{D_t}\sum_{j=1}^{D_t}
    \max\!\bigl(0,\; \gamma - \sqrt{\Var(Z_s^{(j)}) + \delta}\bigr),
\label{eq:var_loss}
\end{equation}
where $Z_s^{(j)}$ denotes the $j$-th column of $Z_s$, $\gamma > 0$ is a
target standard deviation (typically $\gamma = 1$), and $\delta > 0$ prevents
gradient vanishing near zero.  This hinge loss becomes active when any
dimension collapses, pushing its variance back above $\gamma$.

\paragraph{Covariance term $\calL_c$ — enforce decorrelation.}
\begin{equation}
    \calL_c(Z_s) \;=\; \frac{1}{D_t}\, \mathcal{D}(Z_s)
    \;=\; \frac{1}{D_t}\sum_{i \ne j} \hat{C}_{ij}^2,
\label{eq:cov_loss}
\end{equation}
where $\hat{C}$ is the sample covariance of $Z_s$.  This directly minimises the
decorrelation penalty of Definition~\ref{def:decor}.  Unlike the variance term,
$\calL_c$ is active at all times and enforces linear disentanglement as an
ongoing structural constraint.

\paragraph{Whitened alignment term $\calL_a$ — semantic transfer in the disentangled basis.}
\begin{equation}
    \calL_a(Z_s, Z_t) \;=\; \frac{1}{N}\sum_{i=1}^N
    \bigl(2 - 2\,\langle \hat{z}_{s,i}^w,\, \hat{z}_{t,i}^w\rangle\bigr),
\label{eq:align_loss}
\end{equation}
where $\hat{z}^w = z^w / \norm{z^w}$, and $z^w$ denotes the ZCA-whitened,
centred feature vector.  The factor $2 - 2\cos(\cdot) \in [0,4]$ equals the
squared $\ell_2$ distance on the unit sphere and provides stable gradients near
convergence.

\begin{remark}
The alignment term $\calL_a$ differs from plain cosine distillation in two ways:
(1) both student and teacher are whitened using their respective batch covariances
before comparison, and (2) the whitening equalises the gradient magnitude across
all feature dimensions (see Section~\ref{sec:gradient_analysis}).
\end{remark}

\subsection{Gradient Analysis: Why Whitened Alignment Matters}
\label{sec:gradient_analysis}

Consider the gradient of the plain cosine alignment loss with respect to the
student feature $z_s$ (before any projection head):
\[
    \nabla_{z_s} \calL_{\cos}
    = -\frac{1}{\norm{z_s}} \Bigl(
        \hat{z}_t - \langle \hat{z}_s, \hat{z}_t\rangle \hat{z}_s
      \Bigr).
\]
This gradient is proportional to the rejection of $\hat{z}_t$ from
$\hat{z}_s$.  In the original (non-whitened) space, dimensions with large
variance in the teacher produce large gradient components, while low-variance
dimensions are effectively ignored.

After ZCA whitening, the teacher features satisfy $\Cov(Z_t^w) = I$, so all
dimensions have equal variance.  The gradient of the whitened alignment loss
\[
    \nabla_{z_s} \calL_a
    = W_s^\top \bigl(\hat{z}_{t}^w - \langle \hat{z}_s^w, \hat{z}_t^w\rangle
      \hat{z}_s^w\bigr),
\]
where $W_s$ is the student's batch whitening matrix, has magnitude spread
uniformly across dimensions.  This ensures the student must match \emph{all}
semantic directions in the teacher, not just the dominant principal components.

\subsection{Interaction Between the Three Terms}
\label{sec:interaction}

The three terms are complementary and form a self-consistent system:

\begin{itemize}
    \item $\calL_c$ drives $C_s \to \diag(\sigma_1, \ldots, \sigma_D)$
          (diagonal covariance).
    \item $\calL_v$ drives $\sigma_j \ge \gamma$ for all $j$
          (no dimension is dead).
    \item $\calL_a$ drives $Z_s^w \to Z_t^w$ (semantic alignment after whitening).
\end{itemize}

Together, $\calL_v + \calL_c$ enforce $C_s \approx \sigma^2 I$, i.e.\ the
student features become linearly disentangled \emph{as a structural property},
independently of any teacher.  The alignment term $\calL_a$ then provides a
semantic \emph{target} in the already-disentangled space.

\begin{proposition}[Convergence Condition]
Assume $\lambda_v, \lambda_c, \lambda_a > 0$ and that the student has sufficient
capacity.  At a stationary point of $\calL_{\mathrm{LD}}$:
\begin{enumerate}
    \item $C_s \approx \sigma^2 I$ (from $\calL_v + \calL_c$).
    \item The student's whitened patch features are angularly aligned with the
          teacher's whitened patch features (from $\calL_a$).
    \item The student's isotropy score satisfies $\rho(Z_s) \approx 1$
          (Definition~\ref{def:pr}).
\end{enumerate}
\end{proposition}

\subsection{Integration with Existing Distillation Objectives}
\label{sec:integration}

In a multi-teacher distillation framework, the total loss combines the LD
objective with standard token-level and patch-level distillation:

\begin{equation}
    \calL_{\mathrm{total}}
    = \underbrace{\calL_{\cos}^{\mathrm{cls}} + \calL_{\mathrm{rel}} + \calL_{\cos}^{\mathrm{patch}}}_{\text{standard feature distillation}}
    \;+\; \underbrace{\calL_{\mathrm{LD}}}_{\text{linear disentanglement}}.
\label{eq:total}
\end{equation}

The LD term is applied selectively to teachers whose features exhibit linear
disentanglement (DINO, DINOv2), while other teachers (CLIP, MAE) are handled by
their standard distillation terms.

Crucially, the projection head $h_\psi$ used in $\calL_{\mathrm{LD}}$ should be
\emph{separate} from the head used in $\calL_{\cos}^{\mathrm{patch}}$.  This is
because:
(i) cosine alignment is invariant to covariance structure and does not penalise
    entangled features;
(ii) having a dedicated LD head allows $\calL_{\mathrm{LD}}$ to shape the
    representation without conflicting with the cosine supervision signal.

%%% ─────────────────────────────────────────────────────────────────────────
\section{Multi-Teacher Linear Disentanglement Distillation}
\label{sec:multi_teacher}
%%% ─────────────────────────────────────────────────────────────────────────

So far we have focused on distilling the linear disentanglement property from a
single DINO teacher.  A natural question arises: can we extend this framework
to \emph{all} teachers, including those (like MAE) whose features are not
naturally linearly separable?  And if so, what is gained by using multiple
linearly-separable targets simultaneously?

\subsection{Information Preservation under ZCA Whitening}
\label{sec:invertibility}

The key observation is that ZCA whitening is an \emph{invertible} linear map.

\begin{proposition}[ZCA Inverse]
\label{prop:zca_inverse}
The inverse of the ZCA whitening matrix
$W_{\mathrm{ZCA}} = U\Sigma^{-1/2}U^\top$ is
\[
    W_{\mathrm{ZCA}}^{-1} \;=\; U\,\Sigma^{1/2} U^\top.
\]
Consequently, the whitened features $Z_w = \tilde{Z}\,W_{\mathrm{ZCA}}$
contain exactly the same information as the original centred features $\tilde{Z}$:
\[
    \tilde{Z} = Z_w\, W_{\mathrm{ZCA}}^{-1}.
\]
\end{proposition}

This establishes that applying ZCA to \emph{any} teacher — not only DINO — is a
lossless operation.  The whitened MAE features $Z_t^{(\mathrm{mae}),w}$, for
example, are in bijective correspondence with the original MAE encoder outputs.
Distilling from $Z_t^{(\mathrm{mae}),w}$ is therefore fully equivalent in
information content to distilling from the raw features.
This exact equivalence holds for \emph{full-rank} ZCA; the rank-$K$ truncation of
Section~\ref{sec:topk} is a deliberate, lossy restriction to the teacher's signal
subspace, trading the noise tail for numerical stability.  Thus ``lossless''
refers to the full-rank map, whereas in practice we retain only the top-$K$
informative directions.

\subsection{Why Whitening is Especially Beneficial for MAE}
\label{sec:mae_benefit}

Although ZCA preserves information for any teacher, the \emph{practical benefit}
of whitening is largest for teachers whose feature covariance is highly
anisotropic — i.e., teachers for which $\sigma_1 \gg \sigma_D$.
MAE is precisely such a teacher.

\paragraph{Spectral analysis of MAE features.}
MAE is trained to reconstruct masked image patches in pixel or frequency space.
Its optimal encoder features tend to concentrate along directions that explain the
most reconstruction variance — low-spatial-frequency components, colour statistics,
and prominent textures.  As a result, the empirical covariance of MAE patch
tokens has a rapidly decaying eigenspectrum:
\[
    \sigma_1^{(\mathrm{mae})} \;\gg\; \sigma_2^{(\mathrm{mae})} \;\gg\; \cdots
    \;\gg\; \sigma_D^{(\mathrm{mae})}.
\]
By contrast, DINOv2's eigenspectrum is considerably flatter, consistent with its
near-isotropic, linearly-separable representation~\citep{oquab2023dinov2}.

\paragraph{Gradient domination in raw-feature distillation.}
Consider an MSE loss between student projection $Z_s$ and MAE teacher features
$Z_t$ in the eigenbasis $U$ of $\Cov(Z_t)$.  The gradient of the loss with
respect to the $k$-th principal component of the student is
\[
    \frac{\partial \calL_{\mathrm{MSE}}}{\partial [Z_s U]_k}
    \;=\; \frac{2}{N}\bigl([Z_s U]_k - [Z_t U]_k\bigr).
\]
The mean squared magnitude of this gradient over the training distribution
is $\mathcal{O}(\sigma_k^{(\mathrm{mae})})$, since the teacher components have
variance $\sigma_k$.  Therefore, the update to the student is dominated by the
$k=1$ direction — the dominant reconstruction pattern — and the remaining $D-1$
directions receive progressively weaker training signal.

\paragraph{Gradient equalisation after ZCA.}
After whitening, $\Cov(Z_t^w) = I$, so every principal component of the whitened
teacher has unit variance.  The gradient magnitudes become uniform:
\begin{equation}
    \E\!\left[\left(\frac{\partial \calL_a}{\partial [Z_s^w]_k}\right)^{\!2}\right]
    \;=\; \mathcal{O}(1) \quad \forall\, k = 1, \ldots, D.
\label{eq:grad_equal}
\end{equation}
This forces the student to fit \emph{all} of MAE's information — both the
dominant reconstruction directions and the fine-grained, higher-frequency
semantic details that raw MSE under-supervises.

\subsection{Adaptive K Selection for Non-Isotropic Teachers}
\label{sec:topk}

Full-rank ZCA whitening, as defined in Section~\ref{sec:zca}, is well-motivated
when the teacher's covariance matrix is well-conditioned (e.g.\ DINOv2).
For teachers like MAE whose eigenspectrum decays rapidly, however, full-rank
whitening causes two problems:
\begin{enumerate}
    \item \textbf{Noise amplification.}  Near-zero eigenvalues $\sigma_k \approx 0$
          produce $\sigma_k^{-1/2} \gg 1$ in the whitening matrix, amplifying
          the directions that carry essentially no signal.
    \item \textbf{Ill-conditioning.}  Even with regularisation $\sigma_k \leftarrow
          \sigma_k + \varepsilon$, forcing the student to match noise directions
          wastes gradient steps and can destabilise training.
\end{enumerate}

We therefore restrict both the whitening and the distillation objective to the
\emph{linearly separable subspace} of the teacher — the span of eigenvectors
whose eigenvalues carry genuine signal.  An adaptive criterion determines where
signal ends and noise begins.

\paragraph{Effective rank as the primary criterion.}
\begin{definition}[Effective Rank]
\label{def:stable_rank}
The \emph{effective rank} of a feature matrix $Z$ with centred covariance
eigenvalues $\sigma_1 \ge \cdots \ge \sigma_D \ge 0$ is
\begin{equation}
    \operatorname{erank}(Z)
    \;=\; \frac{\bigl(\sum_{i=1}^D \sigma_i\bigr)^2}{\sum_{i=1}^D \sigma_i^2}
    \;\in\; [1,\, D].
\label{eq:stable_rank}
\end{equation}
This is the participation-ratio effective rank; it satisfies
$\operatorname{erank}(Z) = r\cdot\rho(Z)$ with $\rho$ and $r=\rank(C)$ as in
Definition~\ref{def:pr}.  We avoid the name ``stable rank'' for this quantity, as
that term conventionally denotes the distinct operator-norm ratio
$\norm{\,\cdot\,}_F^2/\norm{\,\cdot\,}_2^2$.
\end{definition}

The effective rank equals $D$ when all eigenvalues are identical (perfect
isotropy) and equals $r$ when only $r$ eigenvalues are nonzero (rank-$r$ matrix).
It is continuous, differentiable, and bounded, making it a natural proxy for
the effective dimensionality of the feature distribution.

\paragraph{Energy threshold as a secondary criterion.}
The energy-based rank is
\begin{equation}
    K_\text{energy}(\tau)
    \;=\; \min\!\left\{k : \frac{\sum_{i=1}^k \sigma_i}{\sum_{i=1}^D \sigma_i}
    \ge \tau\right\},
\label{eq:k_energy}
\end{equation}
for a threshold $\tau \in (0,1)$ (typically $\tau = 0.90$).  This guarantees
that the selected subspace contains at least $\tau$ of the total feature variance,
providing a signal-energy floor.

\paragraph{Combined adaptive criterion.}
We define the selected dimension as
\begin{equation}
\boxed{
    K \;=\; \operatorname{clip}\!\Bigl(
        \min\bigl(\lfloor\operatorname{erank}(Z_t)\rfloor,\; K_\text{energy}(\tau)\bigr),\;
        K_{\min},\; D
    \Bigr).
}
\label{eq:k_select}
\end{equation}
Taking the \emph{minimum} of the two criteria is conservative: a dimension is
included only if \emph{both} criteria consider it to carry signal.
The floor of the effective rank provides one estimate of the active
dimensionality; the energy threshold provides a complementary upper bound based
on cumulative variance.  Together they exclude the noise tail robustly, even when one
criterion alone might be too permissive.

Table~\ref{tab:k_values} reports typical $K$ values for the two primary patch
teachers.

\begin{table}[h]
\centering
\caption{Typical adaptive K values for different teachers ($D=768$, $\tau=0.90$).
The MAE values are approximate and depend on training dataset and model size.}
\label{tab:k_values}
\begin{tabular}{lcccc}
\toprule
Teacher & $\operatorname{erank}(Z_t)$ & $K_\text{energy}(0.90)$ & Selected $K$ & $K/D$ \\
\midrule
DINOv2-B  & $\approx 700$  & $\approx 680$ & $\approx 680$ & $88\%$ \\
MAE-B     & $\approx 120$  & $\approx 250$ & $\approx 120$ & $16\%$ \\
MAE-L ($D{=}1024$) & $\approx 180$ & $\approx 350$ & $\approx 180$ & $18\%$ \\
\bottomrule
\end{tabular}
\end{table}

\paragraph{Top-K distillation in teacher's eigenbasis.}
Given the selected $K$ and the corresponding top-$K$ eigenvectors
$U_k \in \R^{D \times K}$ of the teacher covariance, we modify the distillation
procedure as follows.

\begin{enumerate}
    \item \textbf{Project teacher:}
          $Z_t^{(k)} = \tilde{Z}_t\, U_k \in \R^{N \times K}$. \\
          Apply ZCA whitening to obtain $Z_t^{w(k)} = Z_t^{(k)}\, W_t$
          where $W_t$ is the $K \times K$ ZCA matrix of $Z_t^{(k)}$.

    \item \textbf{Project student:}
          $Z_s^{(k)} = \tilde{Z}_s\, U_k \in \R^{N \times K}$ (gradients flow). \\
          Apply stop-gradient ZCA: $Z_s^{w(k)} = Z_s^{(k)}\, W_s^{\mathrm{sg}}$
          where $W_s^{\mathrm{sg}}$ is computed from $\operatorname{sg}(Z_s^{(k)})$,
          so it is a constant w.r.t.\ gradient computation but the linear map
          $Z_s^{(k)} \mapsto Z_s^{w(k)}$ is differentiable.

    \item \textbf{Losses (all in $K$-dim space):}
          \begin{align}
              \calL_v &= \textstyle\frac{1}{K}\sum_{j=1}^K \max(0,\,\gamma - \mathrm{std}([Z_s^{(k)}]_j)),
              \label{eq:var_topk} \\
              \calL_c &= \textstyle\frac{1}{K}\sum_{i \ne j} [\hat C_s^{(k)}]_{ij}^2,
              \label{eq:cov_topk} \\
              \calL_a &= 2 - 2\,\textstyle\frac{1}{N}\sum_i
              \langle \hat z_{s,i}^{w(k)},\, \hat z_{t,i}^{w(k)}\rangle.
              \label{eq:align_topk}
          \end{align}
\end{enumerate}

\begin{remark}[Stop-gradient whitening]
The whitening matrix $W_s^{\mathrm{sg}}$ is computed from a stop-gradient copy
of $Z_s^{(k)}$.  This means the student's covariance is treated as a constant
when computing gradients of $\calL_a$, and the gradient of the alignment loss
with respect to $Z_s^{(k)}$ is simply $W_s^{\mathrm{sg}} \cdot \nabla_{\hat z}
\calL_a$.  This avoids differentiating through the eigendecomposition (which is
numerically fragile at repeated eigenvalues) while still allowing gradients to
propagate to the student backbone through $Z_s^{(k)} = \tilde{Z}_s U_k$.
\end{remark}

\begin{remark}[Why project both into teacher's basis, not student's?]
One could instead project each of teacher and student into their \emph{own}
top-K eigenbases and align them separately.  We prefer teacher's basis for two
reasons.  First, it ensures the student is trained to reproduce structure in the
directions that are \emph{meaningful for the teacher} — the informative signal
space.  Second, using separate eigenbases requires either a subspace alignment
step (e.g.\ Procrustes rotation) or a CKA-style invariant loss, both of which
reintroduce the covariance-invariance problem we are trying to solve.
\end{remark}

\subsection{Teacher Suitability Analysis}
\label{sec:suitability}

Not all teachers in a multi-teacher framework benefit equally from LD
distillation.  Table~\ref{tab:suitability} summarises the analysis.

\begin{table}[h]
\centering
\caption{Suitability of each teacher for LD distillation.}
\label{tab:suitability}
\begin{tabular}{lcccc}
\toprule
Teacher & Features used & Token count $N$ & Anisotropy & Apply LD? \\
\midrule
DINOv2  & Patch tokens  & $B \times L$         & Low  (DINO naturally isotropic) & \checkmark \\
MAE     & Patch tokens  & $B \times L$         & High (reconstruction bias)      & \checkmark \\
CLIP    & CLS token only & $B$                 & N/A  (covariance ill-conditioned) & $\times$ \\
SAM     & Patch tokens  & $B \times L$         & Medium (boundary-focused)       & Optional \\
\bottomrule
\end{tabular}
\end{table}

\noindent The critical factor is the token count $N$ available per batch for
estimating the covariance $\hat{C} \in \R^{D \times D}$.  For patch-based
teachers, $N = B \times L$ (e.g.\ $16 \times 196 = 3136 \gg D = 768$), giving a
well-conditioned estimator.  For CLIP, which provides only one global CLS token
per image, $N = B \approx 32\text{–}64$, making the $D \times D$ covariance
matrix severely rank-deficient and the ZCA transform numerically meaningless.
LD distillation should therefore be applied to \textbf{patch-token teachers only}.

\subsection{Multi-Teacher LD Loss}
\label{sec:multi_teacher_loss}

Let $\mathcal{T}_{\mathrm{patch}} \subseteq \mathcal{T}$ denote the subset of
teachers that provide patch tokens (e.g.\ $\mathcal{T}_{\mathrm{patch}} =
\{\mathrm{dinov3},\, \mathrm{mae}\}$).  For each teacher
$k \in \mathcal{T}_{\mathrm{patch}}$, we introduce a dedicated projection head
$h_k^{\mathrm{LD}} : \R^D \to \R^{D_k}$ that is \emph{separate} from the
standard patch-alignment head $h_k$.  The multi-teacher LD loss is
\begin{equation}
    \calL_{\mathrm{LD}}^{\mathrm{total}}
    \;=\; \sum_{k \in \mathcal{T}_{\mathrm{patch}}} w_k\,
    \calL_{\mathrm{LD}}\!\bigl(Z_s^{(k)},\; Z_t^{(k)}\bigr),
\label{eq:multi_ld}
\end{equation}
where $Z_s^{(k)} = h_k^{\mathrm{LD}}(f_\theta(\mathcal{B})) \in \R^{N \times D_k}$
and $Z_t^{(k)} = g_k(\mathcal{B}) \in \R^{N \times D_k}$ are the student and
teacher patch features respectively, flattened over the spatial dimension.
Each $\calL_{\mathrm{LD}}$ applies the same three-term formula
(Equation~\eqref{eq:total_loss}) independently.

The total training objective becomes
\begin{equation}
    \calL
    = \underbrace{\sum_{k \in \mathcal{T}} \calL_k^{\mathrm{std}}}_{\text{standard distillation}}
    \;+\;
    \underbrace{\calL_{\mathrm{LD}}^{\mathrm{total}}}_{\text{multi-teacher LD}},
\label{eq:grand_total}
\end{equation}
where $\calL_k^{\mathrm{std}}$ includes CLS cosine, relational, and patch
alignment losses for teacher $k$.

\paragraph{Why separate LD heads?}
There are two reasons.
First, the standard patch loss (cosine or MSE) is invariant to the covariance
structure and can be satisfied by an anisotropic representation.  Sharing weights
between the standard and LD heads would create conflicting gradient signals.
Second, having separate heads means the LD property is a \emph{structural
constraint on the backbone}: at inference time, the LD heads are discarded and
the backbone itself produces linearly-separable features, directly usable by
linear probes.

\subsection{Structural Effect on the Student Backbone}
\label{sec:backbone_effect}

We now describe informally what Equation~\eqref{eq:multi_ld} implies for the
student backbone $f_\theta$.

At a stationary point of $\calL_{\mathrm{LD}}^{\mathrm{total}}$:
\begin{enumerate}
    \item For each teacher $k$, the student projection $Z_s^{(k)}$ satisfies
          $C_s^{(k)} \approx \sigma^2 I$ (from the variance and covariance terms),
          meaning it is linearly disentangled with respect to teacher $k$'s
          feature space.
    \item $Z_s^{(k)}$ is angularly aligned with the ZCA-whitened teacher
          $Z_t^{(k),w}$ (from the alignment term), meaning the student captures
          all of teacher $k$'s semantic content uniformly across eigendirections.
    \item Since the above holds independently for \emph{all} patch teachers, the
          backbone must produce representations that can simultaneously be
          projected into multiple linearly-disentangled spaces — each aligned
          with a different teacher's content.
\end{enumerate}

This is a significantly stronger constraint than standard multi-teacher
distillation: the student cannot satisfy it by collapsing to a low-dimensional
subspace or by concentrating features along a single teacher's dominant
directions.  The multi-teacher LD objective creates an implicit pressure for the
backbone to learn a \emph{universally isotropic} representation that contains
structured, linearly-separable information across multiple visual modalities
(semantics, reconstruction, boundary detection).

\subsection{Practical Hyperparameter Guidance}
\label{sec:hyperparams}

The loss weights $\{\lambda_v, \lambda_c, \lambda_a\}$ in
Equation~\eqref{eq:total_loss} and per-teacher weights $w_k$ in
Equation~\eqref{eq:multi_ld} require tuning.  We offer the following guidance.

\textbf{Relative scale of $\lambda_v : \lambda_c : \lambda_a$.}
Following VICReg~\citep{bardes2022vicreg}, we recommend $\lambda_v = \lambda_a =
25$ and $\lambda_c = 1$.  The covariance term involves $D^2$ off-diagonal entries
and thus has a much larger raw magnitude; the factor of $1/D$ normalisation in
$\calL_c$ (Equation~\eqref{eq:cov_loss}) compensates for this.

\textbf{Per-teacher weight $w_k$.}
We recommend setting $w_k = 1$ for all teachers initially.  If the MAE LD loss
converges faster than the DINO LD loss (because MAE's anisotropy is easily
corrected by whitening), one may reduce $w_{\mathrm{mae}}$ after warm-up.

\textbf{Whitening $\varepsilon$.}
Use $\varepsilon = 10^{-3}$ for both DINO and MAE.  MAE's near-zero eigenvalues
(dormant dimensions) are more common, so a slightly larger $\varepsilon =
5 \times 10^{-3}$ may improve stability for MAE.

%%% ─────────────────────────────────────────────────────────────────────────
\section{Connection to Existing Methods}
\label{sec:related}
%%% ─────────────────────────────────────────────────────────────────────────

\paragraph{VICReg~\citep{bardes2022vicreg}.}
VICReg applies variance, invariance, and covariance regularisation directly to
\emph{two augmented views} of the same image in a self-supervised setting.
Our $\calL_v$ and $\calL_c$ are identical to VICReg's variance and covariance
terms applied to the student.  The key difference is that we replace the
invariance term (which requires two views) with a whitened alignment term that
leverages a pre-trained teacher — enabling the structural property to be
acquired \emph{via distillation} rather than self-supervised contrastive
learning.

\paragraph{Barlow Twins~\citep{zbontar2021barlow}.}
Barlow Twins minimises $\norm{C_{YY} - I}_F^2$ where $C_{YY}$ is the
normalised cross-correlation matrix between two views.
This simultaneously drives diagonal entries to 1 (variance, our $\calL_v$)
and off-diagonal entries to 0 (decorrelation, our $\calL_c$).
We apply the same spirit but decompose the two terms, allowing independent
weighting.

\paragraph{W-MSE~\citep{ermolov2021whitening}.}
Whitening Mean Squared Error (W-MSE) applies online whitening to feature
patches before computing MSE across augmentations.  Our $\calL_a$ is
conceptually similar but differs in that we whiten the \emph{teacher's}
representation — this transfers semantic structure from a frozen expert rather
than enforcing invariance across augmentations.

\paragraph{CKA-based distillation.}
Linear CKA, defined as
\[
    \text{CKA}(X,Y) = \frac{\norm{X^\top Y}_F^2}{\norm{X^\top X}_F \norm{Y^\top Y}_F},
\]
is invariant to orthogonal transformations and to global isotropic scaling of
either argument.  It is, however, \emph{not} invariant to arbitrary invertible
linear maps — that property belongs to CCA, not CKA — so, contrary to a common
shorthand, CKA is \emph{not} whitening-invariant.
Nevertheless CKA is unsuitable as an objective for our purpose.  As a rotation-
and scale-invariant \emph{similarity index} it references no canonical basis and
no absolute scale, so maximising $\text{CKA}(Z_s, Z_t)$ does not pin the student
covariance to any prescribed structure such as $C_s = \sigma^2 I$, and it cannot
distinguish a decorrelated student from an orthogonal rotation of it that
scrambles the basis.  CKA can therefore transfer relational structure but cannot
\emph{enforce} the disentanglement target, which is intrinsically basis- and
scale-dependent; our LD objective explicitly addresses this gap.

%%% ─────────────────────────────────────────────────────────────────────────
\section{Discussion}
\label{sec:discussion}
%%% ─────────────────────────────────────────────────────────────────────────

\paragraph{Why does DINO naturally produce linearly disentangled features?}
DINO's self-distillation objective, combined with the centering and sharpening
operations in its student/teacher update, implicitly prevents collapse
(addressing $\calL_v$) and encourages diverse, distributed representations
across the full feature space.  The centering operation in particular acts as a
running estimate of feature mean, preventing any single direction from
consistently dominating — a mechanism analogous to our $\calL_c$.

\paragraph{Universality: applying LD loss to other teachers.}
The framework naturally extends to any teacher.  For MAE features, which are
good at dense prediction but not as isotropic as DINO, one could apply
$\calL_{\mathrm{LD}}$ to distill a more isotropic version of MAE features.
This would transfer MAE's spatial predictive power while adding the linear
separability structure of DINO.

\paragraph{The role of the separate LD head.}
Using a dedicated projection head for the LD objective has an additional
benefit: the head itself can be discarded at inference time.
The student backbone $f_\theta$ is trained so that its intermediate features
are linearly disentangled — the LD property is baked into the backbone, not the
head.
This is analogous to how DINO uses a separate projection head for its
self-distillation but evaluates the backbone directly with linear probes.

\paragraph{Batch size sensitivity.}
The batch ZCA estimator (Section~\ref{sec:batch_zca}) requires $N \gg D$ for a
well-conditioned covariance matrix.  In practice, for $D = 768$ with
$B \times L = 16 \times 196 = 3136$ tokens per batch, the covariance is
well-estimated.  However, for very small batches or large feature dimensions,
a momentum-based running covariance estimate may be more stable.

%%% ─────────────────────────────────────────────────────────────────────────
\section{Conclusion}
\label{sec:conclusion}
%%% ─────────────────────────────────────────────────────────────────────────

We have formalised the linear disentanglement property through two
complementary axes — isotropy (participation ratio) and decorrelation
(off-diagonal covariance penalty) — and shown that $C = \sigma^2 I$ is the
ideal covariance structure for linear probe performance.

ZCA whitening provides the canonical projection into the linearly disentangled
space: it is the unique orthogonalisation that minimises the $\ell_2$ distance
from the original features while achieving $\Cov(Z_w) = I$ exactly.  Its
optimality (Theorem~\ref{thm:zca_optimal}) means that teacher semantics are
preserved as closely as possible under the whitening constraint.

The three-term distillation objective $\calL_{\mathrm{LD}}$ uses the variance
and covariance terms to impose the structural property on the student, and the
whitened alignment term to transfer teacher semantics in the disentangled basis.
This design bridges self-supervised structural regularisation (VICReg-style) and
feature-level knowledge distillation (DINO-style) in a single, principled
objective.

\bibliographystyle{plainnat}
\begin{thebibliography}{10}

\bibitem[Bardes et al.(2022)]{bardes2022vicreg}
Bardes, A., Ponce, J., and LeCun, Y.
\newblock VICReg: Variance-invariance-covariance regularization for
  self-supervised learning.
\newblock \emph{ICLR}, 2022.

\bibitem[Caron et al.(2021)]{caron2021emerging}
Caron, M., Touvron, H., Misra, I., J\'{e}gou, H., Mairal, J., Bojanowski, P., and Joulin, A.
\newblock Emerging properties in self-supervised vision transformers.
\newblock \emph{ICCV}, 2021.

\bibitem[Ermolov et al.(2021)]{ermolov2021whitening}
Ermolov, A., Siarohin, A., Sangineto, E., and Sebe, N.
\newblock Whitening for self-supervised representation learning.
\newblock \emph{ICML}, 2021.

\bibitem[Oquab et al.(2023)]{oquab2023dinov2}
Oquab, M., Darcet, T., Moutakanni, T., et al.
\newblock DINOv2: Learning robust visual features without supervision.
\newblock \emph{TMLR}, 2024.

\bibitem[Zbontar et al.(2021)]{zbontar2021barlow}
Zbontar, J., Jing, L., Misra, I., LeCun, Y., and Deny, S.
\newblock Barlow twins: Self-supervised learning via redundancy reduction.
\newblock \emph{ICML}, 2021.

\end{thebibliography}

\end{document}
